\begin{problem}[Inverse image preserves exactness]\label{prob:vi16-p08}
Let $f:X\to Y$.  Using the stalk identity $(f^{-1}\G)_x\cong\G_{f(x)}$, prove that inverse image preserves exact sequences of sheaves of abelian groups.
\end{problem}

\ifdefined\FullProblemDossiers
\paragraph{Why this problem matters.}
This is the cleanest functorial application of the stalkwise exactness criterion.

\paragraph{Useful ideas.}
Do not compute sections; compute stalks.

\paragraph{Guided hints.}
\begin{enumerate}
\item Take the inverse-image sequence.
\item Identify its stalk at $x$.
\item Apply exactness of the original sequence at $f(x)$.
\end{enumerate}

\begin{solution}
Given an exact sequence on $Y$, its inverse-image sequence has at $x\in X$ the stalk sequence obtained by evaluating the original stalk sequence at $f(x)$.  That sequence is exact.  Since every inverse-image stalk sequence is exact, the sheaf sequence on $X$ is exact by the stalkwise criterion.
\end{solution}

\paragraph{What happened mathematically.}
A potentially complicated sheaf construction becomes trivial after passing to stalks.

\paragraph{Extensions.}
Generalize to complexes and homology sheaves.

\paragraph{Related problems.}
VI.16.P7 contrasts direct image.

\paragraph{Provenance.}
New advanced synthesis.
\else
\paragraph{Provenance.} New advanced synthesis.
\fi
